\documentclass{beamer}

\newcommand{\week}{Financial Markets 2 of 8}

\title{Bonds 2: Nominal v Real Interest Rates, Supply and Demand for Bonds}
\subtitle{Mishkin Chapter 5}
\author{Econ 357}
\date{\week}

% Reference the shared preamble
\input{preamble}

% Set the footer -- change 
\setbeamertemplate{footline}{
  \leavevmode%
  \vspace{2ex}
  \hbox{%
    % Left box: Econ 457
    \begin{beamercolorbox}[wd=.4\paperwidth,ht=2.5ex,dp=1ex,left]{author in head/foot}%
      \hspace{1em}Econ 357
    \end{beamercolorbox}%
    % Middle box: Week
    \begin{beamercolorbox}[wd=.2\paperwidth,ht=2.5ex,dp=1ex,center]{date in head/foot}%
      \centering\week
    \end{beamercolorbox}%
    % Right box: Slide numbers
    \begin{beamercolorbox}[wd=.4\paperwidth,ht=2.5ex,dp=1ex,center]{date in head/foot}%
      \hfill\insertframenumber{} 
    \end{beamercolorbox}%
  }%
  \vskip0pt%
}

\begin{document}

\frame{\titlepage}



\begin{frame}
    \frametitle{Outline for Financial Markets}
        \begin{enumerate}
            \item Financial Market Returns
            \item Bonds 1: Discounting, Prices and Yields
            \item \fbox{Bonds 2: Nominal v. Real Interest Rates, Supply and Demand}
            \item Bonds 3: Market for Money
            \item Bonds 4: Risk and Term Structure
            \item Equities 1: Dividend Discount Model
            \item Equities 2: P/E Ratio
            \item Equities 3: Theories of Stock Pricing
            
        \end{enumerate}
\end{frame}


\begin{frame}
    \frametitle{Outline for Today's Lecture}
    \begin{enumerate}
        \item History of Interest Rates
        \item Real v. Nominal Interest Rates
        \item Supply and Demand for Bonds
        \item Analysis: Inflation, Expansion, Recession
        \item Corporate Bonds
        \item Practice Questions
    \end{enumerate}

\end{frame}

\begin{frame}[t]
    \frametitle{1. History of Interest Rates}
    \section{History of Interest Rates}

    \centering
    \includegraphics[width=0.9\textwidth]{Econ357/figures/02_mkts_3_history1.png}

\end{frame}

\begin{frame}[t]
    \frametitle{2. Real v. Nominal Interest Rates}
    \section{Real v. Nominal Interest Rates}

    \textbf{Definitions}:
    \begin{itemize}
        \item Nominal interest rates: interest rates not adjusted for inflation.  The interest paid in current dollars.
        \item Real interest rates: interest rates adjusted for \textit{expected} inflation.
    \end{itemize}

    \textbf{Discussion:}
    \begin{itemize}
        \item Reminder: Similar to the discussion in a previous lecture about real and nominal returns.
        \item Calculating real returns is relatively easy.  You only need to know how prices have changed, which is easily measured by the CPI.
        \item In contrast, it is harder to calculate real interest rates, because it is hard to calculate \textit{expected} inflation.
    \end{itemize}

\end{frame}

\begin{frame}[t]
    \frametitle{2. Real v. Nominal Interest Rates}
    \framesubtitle{Fisher Equation}

    \textbf{Fisher Equation}
    $$ i = r + \pi^e$$
    Where
    \begin{itemize}
        \item $i$ is the nominal interest rate
        \item $r$ is the real interest rate
        \item $\pi^e$ is \textit{expected} inflation.
    \end{itemize}
    \vspace{1em}
    
    We are often interested in the real interest rate, as it may be a better indicator of borrowing costs or of expected returns.   It is easy to rearrange the Fisher equation:
    $$ r = i - \pi^e$$

    \blfootnote{The Fisher equation is an approximation.  See the book for a proof}
\end{frame}

\begin{frame}[t]
    \frametitle{2. Real v. Nominal Interest Rates}
    \framesubtitle{Expected Inflation}

    Expected inflation is tough to calculate!
    \vspace{1em}

    Discussion:
    \begin{itemize}
        \item Surveys of businesses, households, and economists provide some guidance.
        \item However, generally respondents don't know what inflation is going to be, and this is especially true for households.
        \item There are market-based measures of expected inflation.   These can be useful.
        \item Treasury Inflation Protected Securities (TIPS) have a principal amount that is adjusted with inflation.   This means that the coupon on TIPS is the real interest rate.
    \end{itemize}

\end{frame}

\begin{frame}[t]
    \frametitle{2. Real v. Nominal Interest Rates}
    \framesubtitle{Expected Inflation}

    \centering
    \includegraphics[width=0.8\textwidth]{Econ357/figures/02_mkts_3_exp_inf1png.png}

\end{frame}

\begin{frame}[t]
    \frametitle{2. Real v. Nominal Interest Rates}
    \framesubtitle{Expected Inflation}

    \centering
    \includegraphics[width=0.9\textwidth]{Econ357/figures/02_mkts_3_exp_inf2.png}

\end{frame}

\begin{frame}[t]
    \frametitle{3. Supply and Demand for Bonds}
    \section{Supply and Demand for Bonds}

    Supply and Demand Curve for Bonds:
    \vspace{10em}
    
    Notes:
    \begin{itemize}
        \item Normal supply and demand curves, how to remember: "downward sloping demand"
        \item Price is on the y-axis
    \end{itemize}

\end{frame}

\begin{frame}[t]
    \frametitle{3. Supply and Demand for Bonds}

    Factors Affecting \textbf{Demand}
    \begin{itemize}
        \item Wealth
        \item Expected returns, relative to other asset classes
        \begin{itemize}
            \item If you expect yields to change in the future, what does that mean for your expected return on bonds?
            \item If expected inflation changes, how does that change your real interest rate?
        \end{itemize}
        \item Risk
        \begin{itemize}
            \item Consider risk relative to other asset classes.  An increase in risk in other asset classes could lead to a \textit{relative} decrease in bond risk.
            \item The expected volatility of bond prices also affects the riskiness of bonds.
        \end{itemize}
        \item Liquidity, relative to other asset classes
    \end{itemize}

\end{frame}

\begin{frame}[t]
    \frametitle{3. Supply and Demand for Bonds}

    \centering
    \includegraphics[width=\textwidth]{Econ357/figures/02_mkts_3_sd_table.png}

\end{frame}

\begin{frame}[t]
    \frametitle{3. Supply and Demand for Bonds}

    Factors Affecting \textbf{Supply}
    \begin{itemize}
        \item Government deficit
    \end{itemize}

\end{frame}

\begin{frame}[t]
    \frametitle{4. Analysis}
    \section{Analysis}
    \framesubtitle{Example}

    \textbf{Recession}
    \begin{itemize}
        \item Affect on Demand?
            \begin{itemize}
                \item Channels:
                \item Shift of demand curve
            \end{itemize}
        \item Affect on Supply?
                    \begin{itemize}
                \item Channels:
                \item Shift of supply curve
            \end{itemize}
        \item Result?
    \end{itemize}

\end{frame}

\begin{frame}[t]
    \frametitle{4. Analysis}
    \framesubtitle{Example}

    \textbf{Increase in Expected Inflation}
    \begin{itemize}
        \item Affect on Demand?
            \begin{itemize}
                \item Channels:
                \item Shift of demand curve
            \end{itemize}
        \item Affect on Supply?
                    \begin{itemize}
                \item Channels:
                \item Shift of supply curve
            \end{itemize}
        \item Result?
    \end{itemize}

\end{frame}

\begin{frame}[t]
    \frametitle{4. Analysis}
    \framesubtitle{Example}

    \centering
    \includegraphics[width=\textwidth]{Econ357/figures/02_mkts_3_ex_ust.png}

\end{frame}

\begin{frame}[t]
    \frametitle{5. Corporate Bonds}
    \section{Supply and Demand for Corporate Bonds}

    We can analyze corporate bonds using the same model.  Corporate bonds have a few additional considerations:
    \vspace{1em}

    \textbf{Demand}
    \begin{itemize}
        \item The riskiness of corporate bonds rises as default risk rises, and falls as default risk falls.
    \end{itemize}
    \vspace{1em}

    \textbf{Supply}
    \begin{itemize}
        \item Higher expected inflation may lead corporations to borrow more (lower real interest rates).   This may lead to an increase in supply.
    \end{itemize}

\end{frame}

\begin{frame}[t]
    \frametitle{5. Supply and Demand for Corporate Bonds}
    
    \centering
    \includegraphics[width=\textwidth]{Econ357/figures/02_mkts_3_ex_corps.png}

\end{frame}

\begin{frame}[t]
    \frametitle{5. Supply and Demand for Corporate Bonds}
    
    \centering
    \includegraphics[width=0.9\textwidth]{Econ357/figures/02_mkts_3_history2.png}

\end{frame}

\begin{frame}[t]
    \frametitle{5. Supply and Demand for Corporate Bonds}
    
    \centering
    \includegraphics[width=0.9\textwidth]{Econ357/figures/02_mkts_3_history3.png}

\end{frame}

\begin{practiceframe}
    \frametitle{6. Practice}
    \section{Practice}

    \textbf{Topics Covered}
    \begin{itemize}
        \item Total Return – definition and calculation
        \item Discounting future values
        \item Bond characteristics
        \item Bond pricing.   Definition of bond yields.
        \item Real v. nominal yields.  Challenge of identifying “expected inflation”
        \item Bond supply and demand
        \item Analysis of changes in bond supply and demand: expected inflation, expansion, recession
    \end{itemize}
    
\end{practiceframe}

\begin{practiceframe}
    \frametitle{6. Practice}

    \begin{enumerate}
        \item The present value of an expected future payment (BLANK) as the interest rate increases.
        \item An increase in the time to the promised future payment (BLANK) the present value of the payment.
        \item The yield to maturity on a 10-year treasury note (with face value = \$100 and annual coupon rate = 2.625\%) is 3.37\%. The market price of this bond must be (BLANK) \$100.
        \item The (BLANK) interest rate is adjusted for expected changes in the price level.
        \item If the nominal rate of interest is 2 percent, and the expected inflation rate is -10 percent, what is the real rate of interest?
        \item If you expect the inflation rate to be 15 percent next year and a one-year bond has a yield to maturity of 7 percent, then the real interest rate on this bond is

    \end{enumerate}

\end{practiceframe}

\begin{practiceframe}
    \frametitle{6. Practice}

\begin{enumerate}
        \setcounter{enumi}{6}
        \item In the bond market, the bond demanders are the (BLANK) and the bond suppliers are the (BLANK).
        \item The supply and demand curves for bonds have the usual slopes: upward and downward, respectively.  At (BLANK) prices of the bond, everything else equal, the (BLANK) is higher.
        \item If wealth increases, the demand for stocks (BLANK) and that of long-term bonds (BLANK), everything else held constant.
        \item An increase in the expected rate of inflation will (BLANK) the expected real return on nominal bonds, everything else held constant.   This will cause demand for bonds to (BLANK).
        \item When the expected inflation rate increases, the demand for bonds (BLANK), the supply of bonds (BLANK), and the interest rate (BLANK), everything else held constant.
        \item Everything else held constant, during a business cycle expansion, the supply of corporate bonds shifts to the (BLANK) as businesses perceive more profitable investment opportunities, while the demand for corporate bonds shifts to the (BLANK) as a result of the increase in wealth generated by the economic expansion.   What is the consequence for the price of corporate bonds?

    \end{enumerate}




\end{practiceframe}



\end{document}